\section{Proposed Method}

The proposed system performs orthographic nativization of English loanwords
through a pipeline of stages: graphemic tokenization, priority-ordered rewrite
rules, and G2P lookup. The input to the system is the English orthographic
form; the output is the Filipino orthographic form prescribed by the spelling
conventions in~\cite{almario2014}. Each stage operates on the output of the
previous, and together they model the transduction that constitutes
orthographic nativization. The linguistic properties described
in~\cref{sec:filipino_orthography} directly inform the design of each stage.

\Cref{fig:pipeline} illustrates the multi-stage pipeline. Graphemic tokenization segments the input
into functional units, handling English multigraphs. The rewrite cascade then
maps each token to its Filipino equivalent using context-sensitive rules where
needed and invoking G2P for ambiguous vowels. Finally, output normalization
ensures conformance to Filipino phonotactic constraints.

\begin{figure}[t]
      \centering
      \includegraphics[width=\columnwidth]{assets/system}
      \caption{System architecture.}
      \label{fig:pipeline}
\end{figure}

\subsection{Graphemic Tokenization}
Before any mapping can take place, the English input must be segmented into its
functional graphemic units. This step is necessary because English's deep
orthography contains multi-letter sequences that function as single graphemic
units. The system recognizes the following classes of multigraphs:
\begin{itemize}
      \item Digraphs: \emph{ph}, \emph{ch}, \emph{sh}, \emph{th}, \emph{ng}, \emph{kn},
            \emph{wr}, \emph{gh}, \emph{ee}, \emph{oo}
      \item Trigraphs: \emph{igh}, \emph{tch}, \emph{dge}, \emph{que}, \emph{ore}
      \item Tetragraphs: \emph{augh}, \emph{ough}, \emph{eigh}
\end{itemize}

The tokenizer applies a greedy, longest-match strategy. For example, the words
``straight'' and ``phishing'' are tokenized as:
\begin{center}
      \emph{s}~\emph{t}~\emph{r}~\emph{a}~\emph{igh}~\emph{t} \\
      \emph{ph}~\emph{i}~\emph{sh}~\emph{i}~\emph{ng}
\end{center}
This ensures that subsequent mapping rules operate on meaningful graphemic
units rather than on raw characters.

\subsection{Rewrite Cascades}
Following tokenization, each grapheme is mapped to its Filipino orthographic
equivalent through a priority-ordered cascade of rule types. The system
distinguishes between context-free rules, which apply invariant mappings
regardless of surrounding graphemes, and context-sensitive rules, which select
outputs based on orthographic environment. The system evaluates rules in the
following priority order, accepting the first match:

\begin{enumerate}
      \item Context-sensitive orthographic rules
      \item Phonetic resolution rules
      \item Context-free orthographic rules
\end{enumerate}

% This ordering reflects a key design principle: more specific rules take
% priority over more general ones. Context-sensitive rules handle the majority of
% ambiguous cases where neighboring graphemes determine the mapping;
% context-free rules serve as the default fallback for graphemes with invariant
% correspondences.

\subsubsection{Context-Free Rules}
Many English consonant graphemes bear an invariant relationship to their
Filipino nativized form. For such graphemes, nativization reduces to a direct
one-to-one substitution requiring no knowledge of surrounding context. For
example:

\begin{itemize}
      \item \(\mathit{b} \mapsto \mathit{b}\): Both instances in ``basketball'' map
            directly, yielding \emph{basketbol}.
      \item \(\mathit{ph} \mapsto \mathit{f}\): The digraph \emph{ph}
            consistently maps to \emph{f}; ``cellphone'' \(\mapsto\) \emph{selfon}.
      \item \(\mathit{k} \mapsto \mathit{k}\): ``keyboard'' \(\mapsto\) \emph{kibord}.
\end{itemize}

Because these mappings are strictly positional and produce the same output
regardless of surrounding grapheme, they are computed without reference to any
context window. Context-free rules account for the majority of consonant
mappings in the system.

\subsubsection{Context-Sensitive Rules}

Several of the foreign-derived letters added in the 1987 expansion of the
Filipino alphabet have context-dependent phonemic values, and English graphemes
more broadly often have realizations that vary with their orthographic
environment. For these cases, the mapping for a given grapheme is selected by
inspecting adjacent graphemes and matching against a set of ordered conditions.

For instance, ``c'' is nativized depending on its surrounding context, in this
case it is the vowel after it:
\begin{itemize}
      \item \(\mathit{c} \cdot (\mathit{i} \mid \mathit{e}) \mapsto \mathit{s}\):
            When \emph{c} precedes \emph{i} or \emph{e}, it maps to \emph{s};
            ``cinema'' \(\mapsto\) \emph{sinema}.
      \item Otherwise: \(\mathit{c} \mapsto \mathit{k}\); ``computer'' \(\mapsto\)
            \emph{kompyuter}.
\end{itemize}

Similar context-dependent patterns apply to other graphemes:
\begin{itemize}
      \item Doubled consonants: When two identical consonants appear adjacently, the second
            is treated as silent and mapped to no output, preserving only the single
            consonant; ``addict'' \(\mapsto\) \emph{adik}.
      \item Silent \emph{e}: Word-final \emph{e} following a consonant is dropped; ``make''
            \(\mapsto\) \emph{meyk}.
\end{itemize}

More complex patterns require multi-token lookahead. For instance, the sequence
\emph{-ate} at the word boundary maps to \emph{-eyt} rather than a sequence of
individual vowel and consonant mappings, reflecting the holistic treatment of
common English suffixes in Filipino orthographic convention.

\subsubsection{Phonetic Resolution Rules}\label{subsec:phonetic_resolution}
Vowels pose a larger challenge. English's deep orthography means that a single
vowel grapheme can correspond to several distinct phonemes, and the five-vowel
Filipino system requires that these distinctions be collapsed into a small
target set. Furthermore, because English is a stress-timed language, unstressed
vowels frequently reduce to the schwa /\textipa{@}/, introducing a systematic
ambiguity that orthography alone cannot resolve.

To address this, the system invokes a G2P lookup for vowel tokens whose
nativized form cannot be determined from spelling context alone. The G2P module
produces a phonetic transcription of the entire
English word. The relevant vowel phoneme is then extracted by aligning it to
the current token position. The mapping proceeds as follows:
\begin{itemize}
      \item Tense vowels and diphthongs are rendered as two-character Filipino sequences:
            \emph{ey}, \emph{ow}, \emph{ay}, \emph{aw}.
      \item Lax monophthongs are mapped to the nearest Filipino simple vowel.
      \item The schwa in word-final position is realized as \emph{e} or \emph{o} depending
            on the written vowel.
      \item R-colored vowels preserve the written vowel letter, deferring to orthographic
            convention.
\end{itemize}

\subsection{Output Normalization}
After all per-token rules have been applied, the system performs a final
normalization pass over the raw output sequence. This pass corrects three
classes of artifacts that can arise from the independent application of rules
across adjacent tokens:
\begin{itemize}
      \item Contiguous duplicates: Identical adjacent graphemes produced by separate rule
            branches are collapsed to a single instance.
      \item Redundant affricate prefixes: When a stop consonant is immediately followed by
            its own affricate form (\emph{d} followed by \emph{dy}, or \emph{t} followed by
            \emph{ts}), the leading stop is removed, as it is already encoded within the
            affricate.
      \item Phonotactic cluster repair: Consonant clusters that violate the Filipino CV/CVC
            syllable template are repaired through epenthetic vowel insertion.
            Specifically, when the IPA transcription contains a schwa and the word-final
            graphemes form an illegal cluster ending in \emph{l} or \emph{r}, the vowel
            \emph{o} is inserted before the liquid.
\end{itemize}

\subsection{Rule Construction}

All rules were hand-authored from prescriptive guidelines~\cite{almario2014}
and descriptive accounts of Filipino
phonology~\cite{llamzon1966,llamzon1997,tayao2004,magpale2024}.